\begin{abstract}
scMINER reports that its mutual information clustering, MICA, beats standard single-cell pipelines, but its benchmark lets the reference labels pick every method's setting. We now ask one question: does MICA still win when each method must choose its own clustering without labels? On the authors' released inputs we reproduce MICA's published scores on 9 of 10 datasets. Given all genes and label-picked settings, Leiden on principal components beats MICA on PBMC14K and Zeisel, while MICA leads on seven of the other eight. Without labels the ranking can flip, mostly through how many clusters each rule picks. On Zeisel, Leiden at its default resolution splits 7 types into 21 clusters, while MICA stays near 10. For scaling, we shortlist each cell's likely neighbors by correlation and compute the exact MICA distance only for that shortlist. On 13,605 cells this finds as many true neighbors as MICA's own search, checked by brute force, and runs 6.2 times faster.
\end{abstract}

\begin{IEEEkeywords}
Single-cell RNA sequencing, mutual information clustering, label-free model selection, nearest neighbor graphs, cell-type annotation
\end{IEEEkeywords}
